\section{기록 영상의 고정 모델 오프라인 사례 연구}\label{sec:case}
본 사례는 리프터 운용 중 저장된 네 RGB 영상 8,910프레임에 Base와 Base+N3를 각각 한 번 적용하여 코너·자세의 연속 출력을 비교하는 오프라인 평가이다. 촬영 당시 특정 팔레트에 맞춰 학습한 운용 모델과 논문 평가 모델을 구분하고, 두 평가 경로의 전처리·카메라·치수·좌표계를 고정하였다. 시간축 추적을 포함한 팔레트 연구\cite{mohamed2020}와 달리 새로운 추적기나 제어기를 제안하지 않는다. 두 방법은 각각 8,772프레임에서 새 자세를 산출했고, 이전 출력 유지 없이 138프레임이 no-pose였다.

CSV의 운용 모델 추정값이나 제어 명령은 독립 정답이 아니다. 영상 시각과의 대응, 물체 원점·축·단위를 확인한 뒤 동작 구간 구분에 사용하고, 독립 거리·각도 또는 직접 보이는 코너의 참조가 확보된 항목만 정확도로 평가한다. 완료된 no-pose 구간의 최장은 두 방법 모두 2.535초였고 25개 저장 프레임을 포함했다. 센서 기록 공백은 이 값에 포함하지 않았으며, 저장되지 않은 프레임의 상태를 추정하지 않았다.

세션 경계를 제외하고 양쪽 자세가 모두 유효한 인접 8,737쌍에서 Base와 N3의 $|\Delta x|$ 중앙값은 각각 0.094와 0.100cm, $|\Delta z|$ 중앙값은 0.298와 0.320cm, $|\Delta\psi|$ 중앙값은 0.108과 0.114도였다. P90은 $x$와 $z$에서 각각 0.609→0.593cm와 2.150→2.059cm로 낮아졌고 방향각은 0.431→0.436도로 높아져 변화가 혼합되어 있다. 이 값은 이동과 정지를 함께 포함한 인접 출력 변화이며 정지 잡음이나 정확도가 아니다. 사람 코너 참조·예측 객체 대응·확인된 정지 구간·독립 물리 참조가 아직 없어 해당 정확도는 \notrun 으로 유지한다. 따라서 이 사례로 추적 정확도, 포크 삽입 성공률 또는 정렬시간 개선을 주장하지 않는다.
\begin{table}[t]
\centering\footnotesize
\caption{네 기록 영상 8,910프레임의 고정 모델 오프라인 출력.}\label{tab:case}
\begin{tabularx}{\columnwidth}{Ycc}\toprule
항목 & Base & Base+N3 \\\midrule
평가 영상 / 저장 프레임 & 4 / 8,910 & 4 / 8,910 \\
가시 코너 중앙값 / P90 (px) & \notrun/\notrun & \notrun/\notrun \\
새 자세 산출률 (\%) & 98.451 & 98.451 \\
최장 완료 no-pose 구간 (초) & 2.535 & 2.535 \\
인접 $|\Delta x|$ 중앙/P90 (cm) & 0.094 / 0.609 & 0.100 / 0.593 \\
인접 $|\Delta z|$ 중앙/P90 (cm) & 0.298 / 2.150 & 0.320 / 2.059 \\
인접 $|\Delta\psi|$ 중앙/P90 (도) & 0.108 / 0.431 & 0.114 / 0.436 \\
정지 구간 전방 변동 (cm) & \notrun & \notrun \\
정지 구간 횡방향 변동 (cm) & \notrun & \notrun \\
정지 구간 방향각 변동 (도) & \notrun & \notrun \\
전방 거리 참조 오차 (cm) & \notrun & \notrun \\
횡방향 참조 오차 (cm) & \notrun & \notrun \\
방향각 참조 오차 (도) & \notrun & \notrun \\
\bottomrule\end{tabularx}
\tabnote{두 방법의 fresh/held/no-pose는 각각 8,772/0/138이다. 인접 변화는 양쪽 출력이 유효한 같은 8,737쌍의 절댓값이며 이동과 정지를 함께 포함한다. 변동은 확인된 정지 구간의 표준편차를 뜻하므로 사람 확인 전까지 x다. 가시 코너 및 물리 정확도도 사람 참조와 객체 대응 전까지 x로 둔다.}
\end{table}
